\chapter{Articulación jerárquica}

\section{Articulación entre niveles}
\begin{definition}
Se llama articulación jerárquica a la relación disciplinada mediante la cual varios niveles, umbrales y transiciones pueden organizarse en una estructura coherente sin perder dependencia canónica.
\end{definition}

\begin{observation}
La articulación jerárquica no añade primitivas nuevas. Su función es ordenar de manera legible cómo se enlazan piezas ya distinguidas dentro del canon intermedio.
\end{observation}

\section{Condición mínima de articulación}
\begin{definition}
Existe articulación jerárquica mínima cuando puede declararse simultáneamente:
\begin{enumerate}
  \item precedencia entre niveles;
  \item traducibilidad suficiente entre ellos;
  \item conservación de herencia estructural;
  \item diferencia operativa no colapsada.
\end{enumerate}
\end{definition}

\begin{observation}
Sin estas condiciones, la jerarquía corre el riesgo de convertirse en una lista yuxtapuesta y no en una organización disciplinada.
\end{observation}

\section{Articulación y comparabilidad}
\begin{definition}
La articulación jerárquica requiere comparabilidad suficiente entre piezas adyacentes o traducibles, aun cuando no exija equivalencia plena entre todos los niveles del sistema.
\end{definition}

\begin{observation}
Esta cláusula evita exigir homogeneidad total allí donde solo se necesita relación trazable.
\end{observation}

\section{Articulación y distinción}
\begin{definition}
Toda articulación jerárquica válida debe conservar las distinciones operativas que justifican la separación entre niveles, funciones o formulaciones.
\end{definition}

\begin{observation}
Si la articulación borra las diferencias que pretendía ordenar, la jerarquía se vuelve nominal y pierde fuerza disciplinada.
\end{observation}

\section{Condición mínima de fallo}
\begin{definition}
Se considera fallida una articulación jerárquica cuando no puede sostener de forma simultánea precedencia, herencia, comparabilidad suficiente y distinción operativa entre las piezas articuladas.
\end{definition}

\begin{observation}
Esta condición protege al Tomo II frente a mapas jerárquicos inflados o meramente decorativos.
\end{observation}

\section{Cláusula de alcance}
\begin{observation}
La articulación jerárquica fija una gramática mínima de enlace entre niveles del canon intermedio. No constituye todavía un mapa exhaustivo de todas las dimensiones, regímenes o familias del sistema.
\end{observation}

